% ============================================================
%  EXAMEN FINAL — CURSO INTRODUCTORIO DE MATEMÁTICAS
%  1er Año — Duración: 90 minutos
%  Temas: Aritmética (naturales y decimales), Números enteros,
%         Potenciación y Ecuaciones de primer grado en ℤ
%  Autor: Ing. Gabriel Astudillo
%  Nivel: avanzado (acorde a los exámenes por tema)
%  Nota: enunciado limpio. Las soluciones están en la
%  "Clave de Respuestas" al final (versión _scr sin respuestas).
% ============================================================

\documentclass[12pt, letterpaper]{article}

% ── Paquetes ─────────────────────────────────────────────────

\usepackage{mathwizards-guia}
\mwguia{Examen Final — Curso Introductorio de Matemáticas}{1er Año $\cdot$ 90 minutos}

% ── Comandos propios de esta guía ───────────────────
\newcommand{\Z}{\mathbb{Z}}

\begin{document}
% ============================================================

% ── Portada / Título ─────────────────────────────────────────
\mwportada{Examen Final}{Curso Introductorio de Matemáticas -- 1er Año}{Aritmética, Números Enteros, Potenciación y Ecuaciones}

\vspace{4pt}
\begin{cuadroinfo}
  \small
  \textbf{Nombre:} \rule{0.55\linewidth}{0.4pt} \quad \textbf{Fecha:} \rule{0.15\linewidth}{0.4pt} \\[6pt]
  \textbf{Tiempo disponible:} 90 minutos \quad$\bullet$\quad \textbf{Puntaje total:} 100 puntos \\[4pt]
  \textbf{Instrucciones:}
  \begin{enumerate}[leftmargin=*]
    \item Lee cada ejercicio con calma antes de resolverlo.
    \item Muestra \emph{todos} tus pasos: se evalúa cómo llegas a la respuesta, no solo la respuesta.
    \item En las ecuaciones, \textbf{comprueba} siempre tu solución.
    \item Presta atención a la \emph{posición} de la coma decimal y a los \emph{ceros} no significativos.
    \item No se permite el uso de calculadora.
    \item Escribe con letra clara y revisa tu examen antes de entregarlo.
  \end{enumerate}
\end{cuadroinfo}

\vspace{8pt}

% ============================================================
\section{Aritmética con Naturales y Decimales (20 puntos)}
% ============================================================

\begin{enumerate}[leftmargin=*]
  \item \textbf{(5 pts)} Calcula el valor exacto, cuidando la posición de la coma:
        \begin{enumerate}[label=\alph*)]
          \item $400{,}005 + 0{,}04005 + 40{,}05 + 4{,}00005 + 0{,}4$
          \item $99{,}9999 + 0{,}00019 + 909{,}0909 + 0{,}90909$
        \end{enumerate}

  \item \textbf{(5 pts)} Resuelve las sustracciones (préstamos en cadena):
        \begin{enumerate}[label=\alph*)]
          \item $10000{,}00001 - 998{,}99989$
          \item $305{,}1 - 148{,}980706$
        \end{enumerate}

  \item \textbf{(5 pts)} Calcula los productos con ceros no significativos:
        \begin{enumerate}[label=\alph*)]
          \item $0{,}025 \times 0{,}0040$
          \item $99{,}001 \times 0{,}0099$
        \end{enumerate}

  \item \textbf{(5 pts)} Calcula, respetando la jerarquía y las potencias de diez:
        \begin{enumerate}[label=\alph*)]
          \item $\dfrac{360000 \times 0{,}000025}{0{,}00018 \times 50000}$
          \item $0{,}000007 \times 10^7 \div 0{,}0001 \div 1000 \times 0{,}01 \div 0{,}00001$
        \end{enumerate}
\end{enumerate}

% ============================================================
\section{Números Enteros (20 puntos)}
% ============================================================

\begin{enumerate}[leftmargin=*]
  \setcounter{enumi}{4}
  \item \textbf{(5 pts)} Calcula, simplificando primero los signos dobles:
        \begin{enumerate}[label=\alph*)]
          \item $-(-15) + (-8) - (+6) - (-(-4))$
          \item $|{-12} + 5| - |{-8} - 3| + |-(-7)|$
        \end{enumerate}

  \item \textbf{(5 pts)} Calcula:
        \begin{enumerate}[label=\alph*)]
          \item $(-4) \times 6 \div (-3) \times (-2)$
          \item $(-9) \times 8 \div (-6) \div (-4)$
        \end{enumerate}

  \item \textbf{(5 pts)} Resuelve la operación combinada, evaluando de adentro hacia afuera:
        \[
          [-12 \div (-3) - (-4) \times 2] \div (-2) - (-1)
        \]

  \item \textbf{(5 pts)} En la sierra a las 6 de la mañana la temperatura era de $-8\,^{\circ}$C. Subió $15\,^{\circ}$C, bajó $22\,^{\circ}$C, subió $7\,^{\circ}$C y bajó $3\,^{\circ}$C. \textquestiondown Cuál fue la temperatura final?
\end{enumerate}

% ============================================================
\section{Potenciación (15 puntos)}
% ============================================================

\begin{enumerate}[leftmargin=*]
  \setcounter{enumi}{8}
  \item \textbf{(5 pts)} Calcula, distinguiendo el papel del paréntesis:
        \begin{enumerate}[label=\alph*)]
          \item $(-2)^4$ \quad y \quad $-2^4$
          \item $(-3)^3$ \quad y \quad $-3^3$
        \end{enumerate}

  \item \textbf{(5 pts)} Simplifica usando las propiedades de las potencias:
        \begin{enumerate}[label=\alph*)]
          \item $\dfrac{9^8 \times 9^3}{9^6}$
          \item $\big[(2^3)^4\big]^2$
          \item $5^4 \times 5^2 \div 5^3 \times 5^0$
        \end{enumerate}

  \item \textbf{(5 pts)} Simplifica escribiendo todo en la base común:
        \[
          (64^{5x} \div 128^{x})^{5}
        \]
        \textbf{(Ayuda:} $64 = 2^6$ y $128 = 2^7$.)
\end{enumerate}

% ============================================================
\section{Ecuaciones de Primer Grado (25 puntos)}
% ============================================================

\begin{enumerate}[leftmargin=*]
  \setcounter{enumi}{11}
  \item \textbf{(5 pts)} Resuelve y comprueba: $6x - 15 = 2x + 25$

  \item \textbf{(5 pts)} Resuelve y comprueba: $8 - 5x = 2x - 13$

  \item \textbf{(5 pts)} Resuelve juntando primero los términos semejantes: $5x - 9 + 2x = 3x + 15$

  \item \textbf{(5 pts)} Resuelve y determina si la solución pertenece a $\Z$: $9x - 14 = 4x + 21$

  \item \textbf{(5 pts)} Resuelve: $2[4 - 3(x - 2)] = 4(x + 5)$
\end{enumerate}

% ============================================================
\section{Problemas Integradores (20 puntos)}
% ============================================================

\begin{enumerate}[leftmargin=*]
  \setcounter{enumi}{16}
  \item \textbf{(10 pts)} Sofía va a la tienda con $50$ Bs. y compra tres artículos: uno de $12{,}00$ Bs., otro de $17{,}00$ Bs. y un tercero cuyo precio es el triple de la diferencia entre los dos primeros. Además, decide repartir el vuelto entre sus $4$ hermanos en partes iguales.
        \begin{enumerate}[label=\alph*)]
          \item \textquestiondown Cuánto gastó en total (es decir, $12{,}00 + 17{,}00 + 3 \times (17{,}00 - 12{,}00)$)?
          \item \textquestiondown Cuánto recibe de vuelto?
          \item \textquestiondown Cuánto le corresponde a cada hermano?
        \end{enumerate}

  \item \textbf{(10 pts)} «Pensé un número entero, lo multipliqué por $-2$, al resultado le sumé $14$, luego lo multipliqué todo por $3$ y finalmente le resté $9$, obteniendo $(-27)$.»
        \begin{enumerate}[label=\alph*)]
          \item Escribe la ecuación que representa el problema.
          \item Resuélvela paso a paso y comprueba tu respuesta.
          \item \textquestiondown La solución pertenece a $\Z$?
        \end{enumerate}
\end{enumerate}

\vspace{14pt}
\begin{center}
  {\Large\bfseries\color{mwprimario} ¡Mucho éxito!}\\[4pt]
  {\large\color{gray} Respira profundo, lee con calma y confía en todo lo que practicaste.}
\end{center}

% ============================================================
\ifmwclaves
  \section{Clave de Respuestas}
  % ============================================================

  \subsection*{Aritmética con Naturales y Decimales (Ejercicios 1--4)}

  \begin{enumerate}[leftmargin=*, label=\textbf{\arabic*.}]
    \item a) $444{,}4951$ \quad b) $1010{,}00008$
    \item a) $9001{,}00012$ \quad b) $156{,}119294$
    \item a) $0{,}0001$ \quad b) $0{,}9801099$
    \item a) $1$ \quad b) $700000$ \ ($7 \times 10^5$)
  \end{enumerate}

  \subsection*{Números Enteros (Ejercicios 1--4)}

  \begin{enumerate}[leftmargin=*, label=\textbf{\arabic*.}]
    \item a) $-15 + (-8) - 6 - 4 = -33$ \quad b) $|-7| - |-11| + |7| = 7 - 11 + 7 = 3$
    \item a) $(-4) \times 6 = -24$; $-24 \div (-3) = 8$; $8 \times (-2) = -16$
          \quad b) $(-9) \times 8 = -72$; $-72 \div (-6) = 12$; $12 \div (-4) = -3$
    \item $-12 \div (-3) = 4$; $(-4) \times 2 = -8$; $4 - (-8) = 12$; $12 \div (-2) = -6$; $-6 - (-1) = -5$
    \item $-8 + 15 = 7$; $7 - 22 = -15$; $-15 + 7 = -8$; $-8 - 3 = -11\,^{\circ}$C
  \end{enumerate}

  \subsection*{Potenciación (Ejercicios 1--3)}

  \begin{enumerate}[leftmargin=*, label=\textbf{\arabic*.}]
    \item a) $(-2)^4 = 16$ y $-2^4 = -16$ \quad b) $(-3)^3 = -27$ y $-3^3 = -27$
    \item a) $\dfrac{9^{8+3}}{9^6} = 9^{11-6} = 9^{5}$ \quad b) $\big[(2^3)^4\big]^2 = 2^{24}$ \quad c) $5^{4+2-3+0} = 5^3 = 125$
    \item $64 = 2^6$; $128 = 2^7$ \ $\Rightarrow$ $(2^{30x} \div 2^{7x})^5 = (2^{23x})^5 = 2^{115x}$
  \end{enumerate}

  \subsection*{Ecuaciones de Primer Grado (Ejercicios 1--5)}

  \begin{enumerate}[leftmargin=*, label=\textbf{\arabic*.}]
    \item $6x - 2x = 25 + 15 \Rightarrow 4x = 40 \Rightarrow x = 10$ \checkmark
    \item $8 + 13 = 2x + 5x \Rightarrow 21 = 7x \Rightarrow x = 3$ \checkmark
    \item $7x - 9 = 3x + 15 \Rightarrow 4x = 24 \Rightarrow x = 6$
    \item $9x - 4x = 21 + 14 \Rightarrow 5x = 35 \Rightarrow x = 7$ \ ($7 \in \Z$, sí)
    \item $4 - 3(x - 2) = -3x + 10$; $2(-3x + 10) = 4x + 20 \Rightarrow -6x + 20 = 4x + 20 \Rightarrow x = 0$
  \end{enumerate}

  \subsection*{Problemas Integradores (Ejercicios 1--2)}

  \begin{enumerate}[leftmargin=*, label=\textbf{\arabic*.}]
    \item a) $12{,}00 + 17{,}00 = 29{,}00$; diferencia $= 17{,}00 - 12{,}00 = 5{,}00$; tercero $= 3 \times 5{,}00 = 15{,}00$.
          Total: $29{,}00 + 15{,}00 = 44{,}00$ Bs.
          \quad b) Vuelto: $50 - 44{,}00 = 6{,}00$ Bs.
          \quad c) $6{,}00 \div 4 = 1{,}50$ Bs. por cada hermano.
    \item a) $\big[(-2x + 14) \times 3\big] - 9 = -27$
          \quad b) $3(-2x + 14) - 9 = -27 \Rightarrow -6x + 42 - 9 = -27 \Rightarrow -6x + 33 = -27 \Rightarrow -6x = -60 \Rightarrow x = 10$ \checkmark
          \quad c) $10 \in \Z$, sí.
  \end{enumerate}

  \vspace{10pt}
  \begin{cuadrosolucion}
    \textbf{\textquestiondown Cómo te fue?} Si llegaste hasta el final, \textquestiondown felicidades! Recuerda: los errores no son fracasos, son pistas para aprender. Revisa la posición de la coma, la regla de los signos, las propiedades de las potencias y comprueba siempre tus ecuaciones.
  \end{cuadrosolucion}

\fi
\end{document}
